\section{Введение}
\label{sec:intro}

Функциональный симулятор специализированного ускорителя используется
для разработки программного обеспечения до доступа к устройству.
Замена реализации матричного умножения на GPU может сократить время
исполнения симулятора, но меняет порядок операций с плавающей
точкой~\cite{Higham2002,Goldberg1991}. Поэтому производительность такой
замены необходимо оценивать вместе с численным расхождением результата.

Исследуются три реализации GEMM: последовательная на CPU,
cuBLAS~\cite{cuBLAS2024} и специализированное CUDA-ядро~\cite{CUDA2024}.
Они сопоставляются внутри общей цепочки компиляции и с исполнением
модели без XLA. Объект исследования~--- функциональный симулятор
промышленного ускорителя с заменяемым вычислительным блоком;
его закрытая реализация не раскрывается. Методика сравнения задаётся
через входы, выходы и градиенты модулей и не требует описания
внутреннего устройства симулятора.

Локальное расхождение не определяет итоговое отклонение сети:
его перенос зависит от якобианов, направлений ошибок и их взаимной
зависимости. Для разделения этих факторов используются независимый
повтор слоя на опорных тензорах и отдельный запуск с введением
измеренных возмущений.

Работа включает три результата:
\begin{enumerate}\setlength\itemsep{2pt}
    \item формализацию переноса ошибок с явными условиями линейного
        и статистического приближений;
    \item измерительную схему, разделяющую сравнение путей исполнения,
        локальные ошибки и их распространение;
    \item компонентные испытания 14 архитектурных вариантов, измерения
        времени симуляции и коэффициентов распространения.
\end{enumerate}
Исследуются численные выходы и градиенты отдельных проходов.
Сохранение качества на целевой задаче и сходимость многошагового
обучения не проверяются.
